\section{Program Functionality}\label{sec:functionality}

This section walks through the program as the user meets it, one function at a time, with a screenshot of each. The
screenshots of the brief's test cases are referred to again in Section~\ref{sec:testing}.

The program has a single window (Fig.~\ref{fig:app-05-all-tasks}). Its header takes the place of the system's title
bar and carries the program's mark and name, with the system's minimise, maximise and close buttons drawn on it. A
sidebar on the left lists the filters under \emph{Tasks}: All tasks, Pending, In Progress and Completed, each with the
number of tasks it holds, so that the counts per status are always in view. At the sidebar's foot is the list of
keyboard shortcuts. The rest of the window is the list: a heading (``All tasks'', or the status the list is filtered
by) with a line under it saying what the list shows, such as ``4 tasks, newest first.'', the \emph{Refresh} and
\emph{Add task} buttons, the search field (``Search titles and descriptions''), and the tasks, one row each, under the
columns Task, Status and Updated. The row under the pointer, and the chosen row, show an \emph{Edit} and a
\emph{Delete} button at their end. Forms and confirmations open as dialogs over the window, which is dimmed behind
them. A notice the program has to give, such as a task that no longer exists, appears as a red bar over the list,
with a \emph{Dismiss} button.

\subsection{An empty list}

When there are no tasks, the list is replaced by a message: ``No tasks yet'', with the hint ``Add your first with Add
task or Ctrl+N, such as “Submit the OEL report”.'' (Fig.~\ref{fig:app-01-empty-list}). A search or a filter that
matches nothing says so in the same place: ``Nothing matches “report”'' for a search, or ``No pending tasks'' and the
like for an empty status.

\screenshotwide{app-01-empty-list}
  {The window with no tasks: the list shows a message instead of an empty table (test case 4).}
  {Test case 4: the empty list and its message}

\subsection{Adding a task}

The \emph{Add task} button, or Ctrl+N, opens the form for a new task (Fig.~\ref{fig:app-02-add-form}). It has a
\emph{Title} and a \emph{Description} field, each with a count of its characters against its limit, ``12 / 100'' for
example, counted after trimming as the service counts them and turning red past the limit, and the note ``It starts
as Pending; change its status from the list as the work goes.'' The \emph{Add task} button in the form saves it (as do
Enter in the title and Ctrl+Enter anywhere in the form), and \emph{Cancel} or Esc closes the form without saving. The new task appears at
the top of the list, chosen (Fig.~\ref{fig:app-03-task-added}).

\screenshot{app-02-add-form}
  {The form for a new task, filled in for a second task.}
  {The add form filled in}

\screenshotwide{app-03-task-added}
  {The first task, ``Task 1'', added to the list (test case 1).}
  {Test case 1: the new task in the list}

A title left empty, or made only of spaces, is refused by the service, and its message is shown in the form under the
fields: ``A task needs a title, such as "Submit the OEL report".'' (Fig.~\ref{fig:app-04-empty-title}). The form stays
open with what was typed, so the user can correct it. A title over 100 characters or a description over 1000 is
refused the same way, with a message giving the limit and the length typed.

\screenshot{app-04-empty-title}
  {Saving a task with a description but no title: the form shows ``A task needs a title, such as "Submit the OEL
   report".'' (test case 2).}
  {Test case 2: the validation message}

\subsection{Viewing all tasks}

All tasks are listed newest first (Fig.~\ref{fig:app-05-all-tasks}). Each row shows the title, the first line of the
description, the status, when the task was last updated and when it was added, the times written as ``Today, 21:40'',
``Yesterday, 09:05'' or a date. The status is a coloured pill, grey for Pending, blue for In Progress and green
for Completed, and a completed task's title is greyed and struck through.

\screenshotwide{app-05-all-tasks}
  {Three tasks, one in each status; the sidebar counts 1, 1 and 1 (test case 3).}
  {Test case 3: several tasks, all three statuses}

\subsection{Updating a task}

The row's \emph{Edit} button, Enter on the chosen row, a double click or the row's right-click menu opens the edit
form on the task as it is saved now (Fig.~\ref{fig:app-06-edit-form}). It has the same two fields and a \emph{Status}
choice of Pending, In Progress and Completed, and notes when the task was added and last changed. \emph{Save} saves
all three at once
(Fig.~\ref{fig:app-07-task-updated}).

\screenshot{app-06-edit-form}
  {The edit form on an existing task.}
  {The edit form on an existing task}

\screenshotwide{app-07-task-updated}
  {The task with its changed details in the list (test case 5).}
  {Test case 5: the changed details in the list}

If the task has been deleted in a second copy of the program while the form was open, saving closes the form, shows
the service's message over the list, ``This task no longer exists: it may have been deleted in another copy of the
Task Manager.'', and reads the list again, which no longer has the task (Fig.~\ref{fig:app-08-update-missing}). The
same happens when such a task is opened for editing, deleted or given a new status.

\screenshotwide{app-08-update-missing}
  {Saving a change to a task that a second copy of the program has deleted: the red notice says the task no longer
   exists, and the list, read again, shows the two tasks that remain (test cases 6 and 7).}
  {Test case 6: the message for a task that no longer exists}

\subsection{Deleting a task}

The row's \emph{Delete} button, the Delete key on the chosen row, or \emph{Delete task} in the edit form asks first:
a dialog titled ``Delete this task?'' says ``“Submit the OEL report” will be deleted for good. This can't be
undone.'', naming the task, with \emph{Cancel} and \emph{Delete task} buttons (Fig.~\ref{fig:app-09-delete-confirm}). Only \emph{Delete task} deletes it, and the task then disappears from the list, as Fig.~\ref{fig:app-08-update-missing}
shows for the task deleted in the second copy.

\screenshot{app-09-delete-confirm}
  {The confirmation before ``Task 4'' is deleted.}
  {The delete confirmation}

\subsection{Changing the status from the list}

The status in each row is a button: it opens a menu of the three statuses with the task's own ticked, and choosing
another moves the task to it at once, without opening the form. The counts in the sidebar follow, and the filter for
the new status lists the task (Fig.~\ref{fig:app-11-status-change}).

\screenshotwide{app-11-status-change}
  {The In Progress filter chosen in the sidebar, listing the task whose status was changed from the list; the counts
   are 1, 1 and 1.}
  {A status changed from the list, counts updated}

\subsection{Searching and filtering}

Typing in the search field narrows the list, as each character is typed, to the tasks whose title or description
holds the text, whatever its case (Fig.~\ref{fig:app-12-search-filter}); Ctrl+F goes to the field, Esc clears it, and
the down arrow moves into the list. Choosing a filter in the sidebar shows
one status only, and the two combine: the line under the heading then says, for example, ``2 of 5 pending tasks
match “report”.''

\screenshotwide{app-12-search-filter}
  {A search for ``3'' narrowing the list to one of the three tasks.}
  {Search and status filter narrowing the list}

\subsection{Keyboard shortcuts}

Table~\ref{tab:shortcuts} lists the keyboard shortcuts. The sidebar shows the five that work in the window: add,
edit, delete, search and refresh.

\begin{table}[!t]
  \caption{Keyboard shortcuts}\label{tab:shortcuts}
  \centering\footnotesize
  \begin{tabularx}{\linewidth}{@{}YL{0.2\linewidth}L{0.36\linewidth}@{}}
    \toprule
    Action & Key & Where \\
    \midrule
    Add a task & Ctrl+N & Anywhere in the window \\
    Edit the chosen task & Enter & In the list \\
    Delete the chosen task & Delete & In the list \\
    Go to the search & Ctrl+F & Anywhere in the window \\
    Clear the search & Esc & In the search field \\
    Read the list again & F5 & Anywhere in the window \\
    Save the form & Ctrl+Enter & Anywhere in a form \\
    Save the form & Enter & In the title field \\
    Close a form without saving & Esc & In a form \\
    \bottomrule
  \end{tabularx}
\end{table}

\subsection{Error handling}

The user meets three kinds of error, and none of them as a stack trace. A change the service refuses shows the
service's own message: in the form, under its fields, for a form, or over the list for a change made from the list. A
task that no longer exists is reported over the list, which is then read again. Any other failure, which nobody
expects, such as the database not answering, is written to the log and shown over the list as one plain sentence:
``Something went wrong, and nothing was changed. Try again, or restart the Task Manager if it keeps happening.'' The
same holds for an exception nobody caught on the window's thread, which the program hands to the window instead of
letting it print.
